\title{DeepSeekMath: Pushing the Limits of Mathematical Reasoning in Open Language Models}

\begin{abstract}
Mathematical reasoning remains a formidable obstacle for language models owing to its inherent complexity and structured requirements. Here, we present DeepSeekMath~7B, developed by extending the pre-training of DeepSeek-Coder-Base-v1.5 7B using 120B math-focused tokens extracted from Common Crawl, along with a mix of natural language and code corpora. Notably, DeepSeekMath~7B attains a notable 51.7\% on the MATH benchmark at the competition level without the aid of external toolkits or voting mechanisms, closing the gap with leading models such as Gemini-Ultra and GPT-4. When employing self-consistency sampling over 64 generations, DeepSeekMath~7B reaches a score of 60.9\% on MATH. The strong mathematical reasoning of DeepSeekMath~can be credited to two primary contributors: (1) leveraging the extensive potential in web-scale public datasets through a rigorously designed data selection strategy, and (2) introducing Group Relative Policy Optimization (GRPO), an enhancement over Proximal Policy Optimization (PPO) that not only boosts mathematical reasoning, but also improves memory efficiency over traditional PPO.
\end{abstract}

\section{Introduction}

The emergence of large language models (LLM) has dramatically transformed the landscape of mathematical reasoning within artificial intelligence, driving considerable progress across benchmarks such as quantitative reasoning~\citep{MATH} and geometry reasoning~\citep{trinh2024solving}. Beyond this, LLMs are increasingly vital as tools for supporting humans in tackling intricate mathematical tasks~\citep{tao}. Nonetheless, prominent systems like GPT-4~\citep{gpt4} and Gemini-Ultra~\citep{gemini} remain inaccessible to the public, and available open-source alternatives continue to lag significantly behind in their performance on these mathematical challenges.

In this paper, we present DeepSeekMath, a specialized language model designed for mathematical domains that demonstrates a substantial advancement over existing open-source models and nearly matches GPT-4’s performance across academic benchmarks. To facilitate this, we assemble the DeepSeekMath Corpus—a comprehensive, high-quality pre-training dataset composed of 120B mathematical tokens. The corpus is curated from Common Crawl (CC) through a fastText-based classification system \citep{joulin2016fasttext}. For the initial round, we train the classifier using positive samples sourced from OpenWebMath \citep{openwebmath}, complemented by a varied set of unrelated webpages as negative examples. We then use the trained classifier to extract more positive samples from the CC; these candidates undergo further refinement via human annotation. This refined dataset is subsequently used to retrain the classifier for enhanced accuracy. Our evaluations show that this extensive corpus maintains high quality, with DeepSeekMath-Base 7B attaining 64.2\% accuracy on GSM8K \citep{gsm8k} and 36.2\% on the MATH competition dataset \citep{MATH}, outperforming the Minerva 540B model \citep{minerva}. The DeepSeekMath Corpus also supports multiple languages, and we observe notable gains on Chinese mathematical benchmarks \citep{wei2023cmath,agieval}. We propose that our methodology for mathematical data curation offers valuable groundwork for further investigations, and anticipate substantial potential for future improvements.

DeepSeekMath-Base is built upon DeepSeek-Coder-Base-v1.5 7B \citep{deepseek-coder}, motivated by the observation that initiating training from a code-focused model yields superior results relative to initializing from a general large language model. Additionally, our findings suggest that mathematical pre-training not only strengthens the model’s proficiency in mathematics but also leads to performance gains on general reasoning tasks, as reflected by improved results on the MMLU \citep{mmlu} and BBH \citep{bbh} benchmarks.

Following this pre-training stage, we enhance DeepSeekMath-Base through mathematical instruction tuning using datasets that include chain-of-thought \citep{cot}, program-of-thought \citep{pot,pal}, and tool-assisted reasoning \citep{tora} samples. The final product, DeepSeekMath-Instruct 7B, surpasses all other models of comparable size (7B) and rivals instruction-tuned models at the 70B parameter scale within the open-source landscape.

We also propose Group Relative Policy Optimization (GRPO), an alternative reinforcement learning approach inspired by Proximal Policy Optimization (PPO) \citep{schulman2017proximal}. GRPO replaces the standard critic with a mechanism that estimates the baseline from scores aggregated over groups, thereby greatly reducing computational demands during training. Utilizing a selected subset of English instruction tuning data, GRPO achieves marked gains over the robust DeepSeekMath-Instruct model, with notable improvements observed in both in-domain tasks (GSM8K: 82.9\% $\rightarrow$ 88.2\%, MATH: 46.8\% $\rightarrow$ 51.7\%) and in out-of-domain settings (e.g., CMATH: 84.6\% $\rightarrow$ 88.8\%) throughout the RL process. Furthermore, we present a cohesive framework that elucidates the connections among various approaches, including Rejection Sampling Fine-Tuning (RFT) \citep{yuan2023scaling}, Direct Preference Optimization (DPO) \citep{dpo}, PPO, and GRPO. This framework reveals that all such strategies can be categorized as direct or streamlined forms of reinforcement learning. Through a range of systematic studies—comparing online and offline paradigms, assessing the effects of supervising outcomes versus reasoning processes, and evaluating single-step against iterative RL—we thoroughly dissect the foundational aspects of this approach. Finally, we articulate the reasons behind the significant improvements gained by integrating RL with instruction tuning, and suggest future pathways for refining RL methodologies based on the insights provided by our unified framework.

\subsection{Contributions}
This work presents advancements in scalable math pre-training and investigates reinforcement learning through empirical analysis.

\textbf{Math Pre-Training at Scale}

\begin{itemize}[topsep=0pt]
    \item We demonstrate that Common Crawl, a publicly available dataset, offers substantial value for mathematical tasks. Through a carefully constructed data filtering process, we assembled the DeepSeekMath Corpus—a curated collection comprising 120B tokens drawn from mathematically relevant web content. This dataset is approximately 7 times larger than the math-oriented subset employed by Minerva \citep{minerva}, and about 9 times the volume of OpenWebMath \citep{openwebmath}.

    \item Our DeepSeekMath-Base 7B model, pre-trained from this dataset, matches the mathematical reasoning capabilities of Minerva 540B \citep{minerva}. This outcome suggests that scaling model parameters alone does not dominate performance in mathematical reasoning; instead, the quality of pre-training data significantly impacts results, allowing more compact models to perform effectively.

    \item We report insights gleaned from training models on mathematical content. Notably, pre-training with code prior to mathematical datasets enhances the model’s problem-solving skills, irrespective of external tool assistance. This observation offers evidence towards the hypothesis: \textit{does code training improve reasoning abilities?}—supporting that it does, at least in the mathematical context.

    \item Despite the frequent use of arXiv papers for model training in mathematics-related research, our experiments find no significant benefit on any mathematical evaluation used in this study.
\end{itemize}

\textbf{Exploration and Analysis of Reinforcement Learning}

\begin{itemize}[topsep=0pt]
    \item This work presents Group Relative Policy Optimization (GRPO), a reinforcement learning approach that operates without a critic model. GRPO leverages group scores to estimate baselines, which leads to a substantial reduction in training resources compared to the commonly used Proximal Policy Optimization (PPO).
    \item We show that GRPO markedly improves our instruction-tuned model, DeepSeekMath-Instruct, utilizing only the instruction-tuning dataset. Additionally, we detect gains in out-of-domain generalization during the reinforcement learning phase.
    \item We establish a unified analytical framework encompassing methods such as RFT, DPO, PPO, and GRPO. Within this framework, we carry out comprehensive studies—including comparisons of online and offline training regimes, outcome-based versus process-based supervision, and single-turn versus iterative reinforcement learning—to rigorously analyze the key constituents of the paradigm.
    \item Drawing on insights from our unified framework, we delve into factors contributing to the success of reinforcement learning and highlight several promising pathways for advancing reinforcement learning in large language models (LLMs).
\end{itemize}

\subsection{Summary of Evaluations and Metrics}

\begin{itemize}[topsep=0pt]
    \item \textbf{English and Chinese Mathematical Reasoning}: We undertake thorough evaluations of our models on both English and Chinese mathematical reasoning benchmarks, which range from elementary to university-level problem sets. For English datasets, we consider GSM8K \citep{gsm8k}, MATH \citep{MATH}, SAT \citep{llemma}, OCW Courses \citep{minerva}, and MMLU-STEM \citep{mmlu}. The Chinese evaluations make use of MGSM-zh \citep{mgsm}, CMATH \citep{wei2023cmath}, Gaokao-MathCloze \citep{agieval}, and Gaokao-MathQA \citep{agieval}. Our evaluations assess the capacity of the models to produce standalone text-based solutions (without recourse to external tools) and to resolve questions using Python programming.

    On the English evaluation sets, DeepSeekMath-Base demonstrates comparable performance to the proprietary Minerva 540B \citep{minerva}, and significantly outperforms all existing open-source base models, including Mistral 7B \citep{mistral} and Llemma-34B \citep{llemma}, independent of whether they were math pre-trained. Noteworthy, DeepSeekMath-Base also leads on Chinese datasets, which we attribute to our broader pre-training data selection—incorporating high-quality multilingual datasets rather than restricting to English math content as in prior work \citep{minerva,llemma}. Through subsequent mathematical instruction tuning and reinforcement learning, DeepSeekMath-Instruct and DeepSeekMath-RL reach new milestones, achieving, for the first time in open-source, accuracy rates exceeding 50\% on the challenging MATH competition-level dataset.
\end{itemize}

\item \textbf{Formal Mathematics}: We assess DeepSeekMath-Base's capabilities on the informal-to-formal theorem proving benchmark introduced by \citep{dsp_proof}, using the miniF2F dataset \citep{minif2f} in conjunction with the Isabelle proof assistant \citep{isabelle}. The results indicate that DeepSeekMath-Base achieves notable performance in few-shot autoformalization.

\item \textbf{Natural Language Understanding, Reasoning, and Code}: To obtain a thorough assessment of the model’s general linguistic, reasoning, and programming proficiencies, DeepSeekMath-Base is evaluated on several key benchmarks: the Massive Multitask Language Understanding (MMLU) suite \citep{mmlu}, encompassing 57 varied multiple-choice tasks; BIG-Bench Hard (BBH) \citep{bbh}, which features 23 intricate tasks that mainly emphasize multi-step reasoning; as well as the HumanEval \citep{codex} and MBPP \citep{mbpp} tasks, which are standard measures for code generation models. Pre-training with mathematical data enhances both general comprehension and reasoning skills.

\section{Math Pre-Training}

\subsection{Data Collection and Decontamination}
Here, we describe the methodology for assembling the DeepSeekMath Corpus utilizing data from Common Crawl. Figure \ref{fig:data_collect} illustrates a stepwise workflow that organizes the extraction of a comprehensive mathematics dataset from Common Crawl, beginning with a modest, high-quality seed corpus consisting of math-related material. Importantly, this framework can also be applied to domains beyond mathematics, such as programming.

Our initial seed corpus is OpenWebMath \citep{openwebmath}, which aggregates high-quality mathematical texts sourced from the web. We use this foundation to train a fastText classifier \citep{joulin2016fasttext} for the purpose of identifying additional mathematical web content that resembles OpenWebMath. Specifically, 500,000 examples from the seed are chosen at random to serve as positive instances, while 500,000 unrelated Common Crawl pages serve as negative samples. Model training is conducted with an open-source tool\footnote{\url{https://fasttext.cc}} using the following parameters: a vector size of 256, a learning rate of 0.1, a maximum n-gram length of 3, a minimum occurrence threshold of 3, and three training epochs. To shrink the initial volume of Common Crawl data, both URL-based deduplication and near-duplicate detection are performed, ultimately yielding a set of 40B unique HTML pages. The trained fastText model is then deployed to identify additional mathematical pages from this reduced set. These candidates are ranked based on the model's predicted score, with only the highest-ranking entries retained for further processing. The amount of data selected at this stage is validated by pre-training experiments using subsets of the top 40B, 80B, 120B, and 160B tokens. In the initial cycle, we opt to select the top 40B tokens for further steps.

Following the initial round of data acquisition, a substantial portion of mathematical web content remains uncollected. This shortfall is primarily due to the limited diversity present in the positive samples used to train the fastText model. To address this, we supplement the initial seed corpus by incorporating additional mathematical web sources, thereby improving the robustness of the fastText model. The process involves segmenting the entirety of Common Crawl into separate domains, where a domain consists of web pages sharing the same base URL. For every domain, we determine the fraction of web pages captured in the first collection round. Domains in which more than 10\% of pages are collected are designated as math-related (e.g., \url{mathoverflow.net}). Next, we conduct a manual annotation of URLs within these selected domains that are directly associated with mathematical material (such as \url{mathoverflow.net/questions}). Any remaining, previously uncollected web pages from these annotated URLs are incorporated into the seed corpus. This strategy enables us to broaden the set of positive examples, which in turn leads to a more capable fastText model with enhanced recall for mathematical content in subsequent collection rounds. By the conclusion of four collection cycles, our dataset comprises 35.5M math-related web pages, encompassing 120B tokens. Notably, the fourth collection round retrieves only about 2\% additional data compared to the third, prompting the decision to terminate further data gathering.

In order to prevent contamination of evaluation benchmarks, we adopt the approach detailed in \cite{deepseek-coder} to systematically exclude web pages containing questions or answers from both English-language mathematical benchmarks such as GSM8K \citep{gsm8k} and MATH \citep{MATH}, as well as Chinese datasets such as CMATH \citep{wei2023cmath} and AGIEval \citep{agieval}. Our filtration procedure eliminates any passage containing an exact 10-gram match to any segment from the benchmarks from our math training dataset. For shorter benchmark strings containing at least 3 grams but fewer than 10, a strict exact match is also employed to remove affected web pages.

\subsection{Quality Validation of the DeepSeekMath~Corpus}

To assess the comparative quality of the DeepSeekMath~Corpus, we conduct pre-training experiments alongside several recently published math-oriented corpora:

\begin{itemize}[topsep=0pt]
    \item \textbf{MathPile} \citep{mathpile}: An 8.9B-token corpus collected from multiple sources—including textbooks, Wikipedia, ProofWiki, CommonCrawl, StackExchange, and arXiv—with arXiv accounting for over 85\% of the tokens.
    \item \textbf{OpenWebMath} \citep{openwebmath}: Mathematical content mined from CommonCrawl, totaling 13.6B tokens.
    \item \textbf{Proof-Pile-2} \citep{llemma}: A composite corpus of OpenWebMath, AlgebraicStack (10.3B tokens of math code), and arXiv articles (28.0B tokens). For our Proof-Pile-2 experiments, we follow the arXiv:Web:Code data ratio of 2:4:1 as described in \cite{llemma}.
\end{itemize}

\subsubsection{Training Setting}
\label{sec:quality-policy}
We specialize a general-purpose language model containing 1.3B parameters—architecturally aligned with DeepSeek LLMs \citep{deepseek-llm} and referred to as DeepSeek-LLM 1.3B—by fine-tuning it with mathematical data. Each mathematical corpus is used to train a distinct model for a total of 150B tokens. All training is performed using the resource-efficient HAI-LLM framework \citep{haillm}. In line with the DeepSeek LLMs' established methodology, we employ the AdamW optimizer \citep{adamW} configured with $\beta_1=0.9$, $\beta_2=0.95$, and a $\mathrm{weight\_decay}$ of $0.1$. The learning rate schedule follows a multi-step approach: after 2{,}000 warmup iterations, the learning rate peaks, then reduces to 31.6\% of the maximum after 80\% of total training steps, and finally to 10.0\% of the peak following 90\% of training. The maximum learning rate is set to $5.3\times10^{-4}$, while training is performed with a batch size of 4 million tokens and a sequence context of 4K tokens.

\subsubsection{Evaluation Results}

\textbf{The DeepSeekMath~Corpus stands out for its superior quality, extensive multilingual representation, and unparalleled scale.}

\begin{itemize}[topsep=0pt]
    \item \textbf{High-quality}:\\
    We assess downstream abilities across eight mathematical benchmarks using few-shot chain-of-thought prompting \cite{cot}. Results presented in Table~\ref{tab:corpora_comparison} demonstrate that the DeepSeekMath~Corpus confers a distinct performance edge to trained models. Further, Figure~\ref{fig:corpora_comparisons} illustrates that, after 50B training tokens (equivalent to one complete epoch of Proof-Pile-2), models fine-tuned on DeepSeekMath outperform those trained on Proof-Pile-2. This suggests that the average data quality within the DeepSeekMath~Corpus surpasses that of existing corpora.
    \item \textbf{Multilingual}:\\
    The DeepSeekMath~Corpus incorporates material spanning multiple languages, with English and Chinese comprising the bulk of the dataset. As indicated by Table~\ref{tab:corpora_comparison}, training on DeepSeekMath~Corpus leads to enhanced mathematical reasoning in both English and Chinese, while prior corpora, largely dominated by English, either deliver minimal improvement or negatively impact reasoning in Chinese.
    \item \textbf{Large-scale}:\\
    In terms of volume, the DeepSeekMath~Corpus significantly exceeds currently available mathematical corpora. According to Figure~\ref{fig:corpora_comparisons}, DeepSeek-LLM 1.3B models trained on DeepSeekMath display both a more rapid progression and sustained performance increases. By contrast, existing baselines, which are considerably smaller, require multiple epochs of repetition during training and plateau quickly in performance.
\end{itemize}

\subsection{Training and Evaluating DeepSeekMath-Base 7B}

This section presents DeepSeekMath-Base 7B, a foundational model engineered for robust reasoning, with a particular emphasis on mathematical problem-solving. We begin by initializing our model with DeepSeek-Coder-Base-v1.5 7B \citep{deepseek-coder} and proceed to train it on a corpus comprising 500B tokens. The composition of our training data is as follows: 56\% originates from the DeepSeekMath Corpus, 4\% from AlgebraicStack, 10\% from arXiv, 20\% consists of Github code, and the final 10\% comprises natural language resources from the English and Chinese sections of Common Crawl. Except for modifications—namely, assigning the maximum learning rate as $4.2\times10^{-4}$ and employing a batch size of 10 million tokens—our training adheres primarily to the configurations detailed in Section \ref{sec:quality-policy}.

We undertake an in-depth evaluation of DeepSeekMath-Base 7B’s mathematical proficiencies. This encompasses its capability to independently generate complete mathematical solutions, address mathematical problems by leveraging tools, and execute formal theorem verification. Additionally, we survey the general capabilities of the base model, including its aptitude in natural language comprehension, logical reasoning, and coding proficiency.

\paragraph{Mathematical Problem Solving with Step-by-Step Reasoning}

DeepSeekMath-Base is assessed on its ability to solve mathematical tasks using few-shot chain-of-thought prompting \citep{cot} across eight English and Chinese benchmarks. These benchmarks span quantitative reasoning tasks (such as GSM8K \citep{gsm8k}, MATH \citep{MATH}, and CMATH \citep{wei2023cmath}) as well as multiple-choice formats (for instance, MMLU-STEM \citep{mmlu} and Gaokao-MathQA \citep{agieval}). The selected benchmarks represent a broad spectrum of mathematical disciplines, ranging from primary school problems to those at the university level.

As presented in Table \ref{tab:base_math_cot}, DeepSeekMath-Base 7B achieves leading results across all eight benchmarks among open-source base models—this set comprises prominent models such as the general-purpose Mistral 7B \citep{mistral} and the recently introduced, math-focused Llemma 34B \citep{llemma}, trained on Proof-Pile-2 \citep{llemma}. Of particular note, DeepSeekMath-Base attains a performance advantage exceeding 10\% in absolute terms over prior open-source base models on the competitive MATH benchmark, even surpassing Minerva 540B \citep{minerva}, a proprietary base model that is 77 times larger, built upon PaLM \citep{palm}, and subjected to additional mathematical pretraining.

\paragraph{Mathematical Problem Solving with Tool Use} We assess models’ abilities in program-assisted mathematical reasoning using GSM8K and MATH through few-shot program-of-thought prompting \citep{pot,pal}. For each problem, models are instructed to compose Python scripts, optionally leveraging libraries like \textit{math} and \textit{sympy} to handle complex calculations. The program’s execution output is regarded as the solution. Results in Table \ref{tab:base_math_pot_proof} reveal that DeepSeekMath-Base 7B surpasses the previous state-of-the-art, Llemma 34B.

\paragraph{Formal Mathematics}

The automation of formal proof writing plays an important role in ensuring both the rigor and efficiency of mathematical verification, a topic of rising interest. We measure DeepSeekMath-Base 7B’s capability in the informal-to-formal proof generation task from \citep{dsp_proof}, which involves transforming an informal description, a formal version, and an informal proof into a formalized proof. This evaluation is carried out on miniF2F \citep{minif2f}, a benchmark targeting formal proofs at the Olympiad level, where for each instance, we use few-shot prompting to produce an Isabelle formal proof. Following the methodology of \cite{dsp_proof}, we have models create proof sketches which are then refined with Sledgehammer \citep{sledgehammer}, an off-the-shelf automated prover. According to Table \ref{tab:base_math_pot_proof}, DeepSeekMath-Base 7B attains robust results in the domain of automated proof formalization.

\paragraph{Natural Language Understanding, Reasoning, and Code}

We investigate the model’s proficiency in natural language understanding on MMLU \citep{mmlu}, reasoning with BBH \citep{bbh}, and code generation on HumanEval \citep{codex} and MBPP \citep{mbpp}. Table \ref{tab:understanding_reasoning_code} demonstrates that DeepSeekMath-Base 7B yields marked improvements on MMLU and BBH relative to DeepSeek-Coder-Base-v1.5 \citep{deepseek-coder}, indicating that specialized mathematics training enhances capabilities in language comprehension and logical reasoning. Furthermore, the model’s continued training with code tokens allows it to match DeepSeek-Coder-Base-v1.5’s results on both coding tasks. In sum, DeepSeekMath-Base 7B achieves substantially better performance than the generalist Mistral 7B \citep{mistral} on all three benchmarks related to reasoning and programming.

\section{Supervised Fine-Tuning}

\subsection{SFT Data Curation}

We developed an instruction-tuning dataset for mathematics, comprising both English and Chinese problems that span multiple mathematical domains and a range of difficulty levels. Each problem is accompanied by a detailed solution, formatted as chain-of-thought (CoT) \citep{cot}, program-of-thought (PoT) \citep{pot,pal}, or employs tool-integrated reasoning approaches \citep{tora}. In total, the dataset encompasses 776K training instances.

\begin{itemize}[topsep=0pt]
    \item \textbf{English mathematical datasets}: We enhance GSM8K and MATH datasets by providing tool-integrated solution annotations, supplement the dataset with selected examples from MathInstruct \citep{MathInstruct}, and incorporate the training partition of Lila-OOD \citep{lila}, where each problem is solved with either CoT or PoT strategies. Our English dataset covers an extensive set of mathematical disciplines such as algebra, probability, number theory, calculus, and geometry.
    \item \textbf{Chinese mathematical datasets}: We aggregate Chinese K-12 mathematics problems, spanning 76 distinct subfields including areas like linear equations. Each is annotated with solutions that follow both CoT and tool-integrated reasoning methods.
\end{itemize}

\subsection{Training and Evaluating DeepSeekMath-Instruct 7B}

Here, we describe the supervised instruction tuning of DeepSeekMath-Instruct 7B, which is initialized from DeepSeekMath-Base. For training, input samples are randomly concatenated up to a context window of 4K tokens. The model undergoes 500 optimization steps with a batch size of 256, maintaining a fixed learning rate of 5e-5.

The model’s mathematical abilities, with and without tool usage, are assessed using four quantitative reasoning benchmarks in both English and Chinese. For comparison, we evaluate DeepSeekMath-Instruct 7B alongside the most advanced models available:

\begin{itemize}[topsep=0pt]
    \item \textbf{Closed-source models}: We consider several leading systems: (1) members of the GPT family, notably GPT-4 \citep{gpt4} and the GPT-4 Code Interpreter~\footnote{\url{https://openai.com/blog/chatgpt-plugins\#\#code-interpreter}}, (2) Gemini Ultra and Pro \citep{gemini}, (3) Inflection-2 \citep{inflection-2}, (4) Grok-1~\footnote{\url{https://x.ai/model-card}}, and (5) Baichuan-3~\footnote{\url{https://www.baichuan-ai.com}}, alongside (6) the most recent GLM-4~\footnote{\url{https://open.bigmodel.cn/dev/api\#glm-4}}
    
    from the GLM family \citep{glm}.
\end{itemize}

Most of these models are designed for broad applications and have typically been aligned using various alignment methods.

\item \textbf{Open-source models} encompass: foundational general models such as (1) DeepSeek-LLM-Chat 67B \citep{deepseek-llm}, (2) Qwen 72B \citep{qwen}, (3) SeaLLM-v2 7B \citep{seallm}, and (4) ChatGLM3 6B \citep{chatglm3}; in addition to a set of models specifically adapted for mathematics: (5) InternLM2-Math 20B~\footnote{\url{https://github.com/InternLM/InternLM-Math}}, derived from InternLM2 with targeted math training followed by instruction tuning, (6) Math-Shepherd-Mistral 7B, which incorporates PPO training \citep{schulman2017proximal} into Mistral 7B \citep{mistral} and leverages a process-supervised reward model, (7) the WizardMath family \citep{wizardmath}, which augments mathematical reasoning in Mistral 7B and Llama-2 70B \citep{llama2} using evolve-instruct (a form of instruction tuning utilizing AI-generated instructions) and PPO training with primary datasets from GSM8K and MATH, (8) MetaMath 70B \citep{metamath}, created by fine-tuning Llama-2 70B on an expanded GSM8K and MATH corpus, (9) ToRA 34B \cite{tora}, which is a CodeLlama 34B variant further fine-tuned for mathematical reasoning with tool integration, and (10) MAmmoTH 70B \citep{MathInstruct}, which applies MathInstruct-based instruction tuning to Llama-2 70B.

As detailed in Table \ref{tab:sft_rl_math}, within an evaluation setting that prohibits the use of external tools, DeepSeekMath-Instruct 7B delivers robust stepwise reasoning capabilities. Importantly, on the competition-grade MATH benchmark, our model exceeds the performance of every other open-source model and the vast majority of proprietary counterparts (including Inflection-2 and Gemini Pro) by a margin of at least 9 percentage points. This holds even against much larger models (for instance, Qwen 72B) or those leveraging advanced reinforcement learning focused on mathematics (like WizardMath-v1.1 7B). Although DeepSeekMath-Instruct approaches the performance levels of proprietary Chinese models such as GLM-4 and Baichuan-3 on MATH, it still lags behind the results achieved by GPT-4 and Gemini Ultra.

When evaluated with the permission for combined natural language reasoning and code-based tools, DeepSeekMath-Instruct 7B nearly attains 60\% accuracy on MATH, outperforming all previously released open-source models. On other datasets, our model's results are on par with DeepSeek-LLM-Chat 67B, the preceding open-source state-of-the-art, despite being one-tenth its size.

\section{Reinforcement Learning}

\subsection{Group Relative Policy Optimization} Reinforcement learning (RL) has consistently demonstrated its utility in enhancing the mathematical reasoning abilities of LLMs after they have completed the Supervised Fine-Tuning (SFT) phase \citep{wang2023math,wizardmath}. Here, we present Group Relative Policy Optimization (GRPO), a resource-efficient and effective RL approach.

\subsubsection{From PPO to GRPO} Proximal Policy Optimization (PPO) \citep{schulman2017proximal} is a widely adopted actor-critic RL algorithm utilized during the fine-tuning process for LLMs \citep{ouyang2022training}. Specifically, it seeks to improve model parameters by maximizing the following surrogate objective:

\begin{equation}
\footnotesize
    \mathcal{J}_{PPO}(\theta) = \mathbb{E}{[q \sim P(Q), o \sim \pi_{\theta_{old}}(O|q)]} \frac{1}{|o|} \sum_{t=1}^{|o|} \min \left[ \frac{\pi_\theta(o_{t} | q, o_{<t})}{\pi_{\theta_{old}}(o_{t} | q, o_{<t})} A_{t}, \text{clip} \left( \frac{\pi_\theta(o_{t} | q, o_{<t})}{\pi_{\theta_{old}}(o_{t} | q, o_{<t})}, 1 - \vare

Here, $\pi_{\theta}$ denotes the current policy model and $\pi_{\theta_{old}}$ refers to the previous policy. The variables $q$ and $o$ represent, respectively, the questions and sampled responses drawn from the question set and from the old policy $\pi_{\theta_{old}}$. The parameter $\varepsilon$ serves as a hyper-parameter related to clipping, originally introduced in PPO to enhance training robustness. The term $A_t$ stands for the advantage, which is estimated using Generalized Advantage Estimation (GAE) \citep{gae} from future rewards $\{r_{\ge t}\}$ along with a learned value function $V_{\psi}$. As a result, the PPO framework necessitates joint training of both the policy and a value function. To counteract the risk of the reward model being excessively optimized, a widely adopted practice involves incorporating a per-token KL penalty relative to a reference model into the reward for each token, following \citep{ouyang2022training}:

\begin{equation}
    r_{t} = r_\varphi(q, o_{\le t}) - \beta \log\frac{\pi_{\theta}(o_{t}|q, o_{<t})}{\pi_{ref}(o_{t}|q, o_{<t})},
\label{eq:PPO-reward}
\end{equation}

where $r_\varphi$ specifies the reward model, $\pi_{ref}$ indicates the reference model (commonly initialized as the SFT model), and $\beta$ controls the strength of the KL penalty.

Because the PPO value function typically matches the policy model in scale, this leads to considerable demands on memory and computation. Additionally, during reinforcement learning, the value function operates as a baseline to minimize variance in advantage estimation. However, for LLMs, the reward model usually issues a reward only for the final token, which poses difficulties for learning a value function that is accurate at every time step. To resolve this challenge, we introduce Group Relative Policy Optimization (GRPO) (see Figure \ref{fig:grpo}), which dispenses with the need for a separate value function. Instead, GRPO leverages the mean reward across a collection of sampled outputs, generated in response to the same question, to serve as a baseline. Specifically, for any given question $q$, GRPO draws a group of outputs $\{o_1, o_2, \cdots, o_G\}$ from $\pi_{\theta_{old}}$, and then updates the policy by maximizing the following objective:

\begin{equation}
\footnotesize
\begin{split}
    \mathcal{J}_{GRPO}(\theta) &= \mathbb{E}{[q \sim P(Q), \{o_i\}_{i=1}^G \sim \pi_{\theta_{old}}(O|q)]}  \\
    & \frac{1}{G}\sum_{i=1}^G\frac{1}{|o_i|} \sum_{t=1}^{|o_i|} \left\{ \min \left[ \frac{\pi_\theta(o_{i,t} | q, o_{i,<t})}{\pi_{\theta_{old}}(o_{i,t} | q, o_{i,<t})} \hat{A}_{i,t}, \text{clip} \left( \frac{\pi_\theta(o_{i,t} | q, o_{i,<t})}{\pi_{\theta_{old}}(o_{i,t} | q, o_{i,<t})}, 1 - \varepsilon, 1 + \varepsilon \right)  \hat{A}_{i,t} \right] - \beta \mathbb{D}_{KL}\left[\pi_{\theta} || \pi_{ref}\right]\right\} ,
\end{split}
\label{eq:GRPO-obj}
\end{equation}

Here, $\varepsilon$ and $\beta$ serve as hyper-parameters. The term $\hat{A}_{i,t}$ refers to the advantage, which in this method is calculated by comparing only rewards among outputs within each group; its computation will be elaborated in later subsections. The GRPO strategy for advantage computation is consistent with the inherent comparative design of reward models, as these are often trained using datasets comprised of output comparisons given the same prompt. Furthermore, instead of embedding the KL term within the reward function, GRPO directly introduces the KL divergence between the current policy and the reference policy as a regularization component in the loss function, thereby sidestepping added complexity in $\hat{A}_{i,t}$'s calculation. Importantly, rather than relying on the KL penalty as in (\ref{eq

\begin{algorithm}[t]
  \small
  \caption{Iterative Group Relative Policy Optimization}
  \textbf{Input} initial policy model $\pi_{\theta_{\text{init}}}$; reward models $r_\varphi$; task prompts $\mathcal{D}$; 
  hyperparameters $\varepsilon$, $\beta$, $\mu$
  \begin{algorithmic}[1]
    \State Initialize policy model $\pi_\theta \leftarrow \pi_{\theta_{\text{init}}}$
    \For{iteration = 1, \dots, I}
      \State Assign the current policy model as the reference: $\pi_{ref} \leftarrow \pi_{\theta}$
      \For{step = 1, \dots, M}
      \State Draw a batch $\mathcal{D}_b$ from $\mathcal{D}$
      \State Set $\pi_{\theta_{old}} \leftarrow \pi_{\theta}$ 
      \State For each prompt $q \in \mathcal{D}_b$, sample $G$ responses $\{o_i\}_{i=1}^G \sim \pi_{\theta_{old}} (\cdot \mid q)$
      \State Use $r_\varphi$ to compute the reward for each output, resulting in $\{r_i\}_{i=1}^{G}$
      \State Calculate $\hat{A}_{i,t}$ for the $t$-th token of each $o_i$ via group relative advantage estimation
      \For{GRPO iteration = 1, \dots, $\mu$}
        \State Update $\pi_\theta$ by optimizing the GRPO objective (Equation \ref{eq:GC-GRPO})
      \EndFor
    \EndFor \State Continue training $r_\varphi$ with a replay mechanism.
    \EndFor 
  \end{algorithmic}
  \textbf{Output} $\pi_\theta$
  \label{alg:iter-grpo}
\end{algorithm}

\subsubsection{Outcome Supervision RL with GRPO} More specifically, given any question $q$, $G$ candidate outputs $\{o_1, o_2, \cdots, o_G\}$ are drawn from the previous policy version $\pi_{\theta_{old}}$. Each response is assigned a reward by the reward model, producing the set $\mathbf{r}=\{r_1, r_2, \cdots, r_G\}$. The obtained rewards are then normalized by subtracting their mean and dividing by the standard deviation within the group. In outcome supervision, each output $o_i$ receives its normalized reward at the end of the sequence, and every token in $o_i$ shares the same advantage value equal to the normalized reward; that is, $\hat{A}_{i, t} = \widetilde{r}_i = \frac{r_i- {\rm mean}(\mathbf{r})}{{\rm std}(\mathbf{r})}$. The policy parameters are subsequently optimized to maximize the criterion defined in equation (\ref{eq:GRPO-obj}).

\subsubsection{Process Supervision RL with GRPO} Since outcome supervision only issues feedback at the completion of each sequence, this may be inadequate for steering the model during sophisticated mathematical reasoning tasks. Building upon the framework of \cite{wang2023math}, we incorporate process supervision, which assigns a reward to each intermediate reasoning step. To formalize, for a prompt $q$ and a set of $G$ sampled completions $\{o_1, o_2, \cdots, o_G\}$, a process-based reward model evaluates each step, resulting in structured rewards $\mathbf{R} = \{ \{r_1^{index(1)},\cdots,r_1^{index(K_1)}\}, \cdots,  \{r_G^{index(1)},\cdots,r_G^{index(K_G)}\} \}$, where $index(j)$ marks the terminal token of step $j$ in output $i$, and $K_i$ represents the step count for $o_i$. All such step-level rewards are normalized via mean and standard deviation, so that $\widetilde{r}_i^{index(j)} = \frac{r_i^{index(j)} - {\rm mean(\mathbf{R})}}{{\rm std(\mathbf{R})}}$. For process supervision, the advantage at position $t$ of output $o_i$ is defined as the sum of normalized rewards for all subsequent steps, $\hat{A}_{i, t} = \sum_{index(j) \ge t} \widetilde{r}_i^{index(j)}$. Finally, policy improvement

\subsubsection{Iterative RL with GRPO} During reinforcement learning, the previously trained reward model may become inadequate for effectively guiding the evolving policy model. To address this, we investigate an iterative RL approach utilizing GRPO. As detailed in Algorithm \ref{alg:iter-grpo}, the iterative GRPO framework involves generating fresh datasets for training the reward model by sampling outputs from the current policy model. The reward model is then incrementally updated, employing a replay strategy that retains 10\% of past training data. Subsequently, the policy model is used as the new reference, and its training continues under supervision from the updated reward model.

\subsection{Training and Evaluating DeepSeekMath-RL}

Our reinforcement learning experiments employ the DeepSeekMath-Instruct 7B model as a foundation. The RL training dataset comprises approximately 144K chain-of-thought-formatted questions pertaining to GSM8K and MATH, extracted from the SFT dataset. Questions outside this domain are omitted to analyze how RL affects benchmark performance when certain types of data are absent during RL training. The procedure for assembling the reward model's training set aligns with \citep{wang2023math}. The reward model is initially trained using DeepSeekMath-Base 7B, adopting a learning rate of 2e-5. For GRPO, the policy model's learning rate is configured as 1e-6, and the KL coefficient is set to 0.04. For every question, we generate $64$ sample outputs. Both the maximum input length and training batch size are fixed at 1024. The policy model undergoes just one update per exploration iteration. DeepSeekMath-RL 7B is assessed on the same set of benchmarks as DeepSeekMath-Instruct 7B. In the case of DeepSeekMath-RL 7B, GSM8K and MATH with chain-of-thought reasoning serve as in-domain evaluations, while the remaining benchmarks are considered out-of-domain.

Table \ref{tab:sft_rl_math} presents a comparative evaluation of open- and closed-source models employing both chain-of-thought and tool-assisted reasoning across English and Chinese datasets. Key observations include: 1) DeepSeekMath-RL 7B achieves accuracy rates of 88.2\% on GSM8K and 51.7\% on MATH using chain-of-thought strategies. These results outperform all other open-source models in the 7B–70B parameter range and surpass most closed-source counterparts as well. 2) Notably, DeepSeekMath-RL 7B's training regimen is restricted solely to chain-of-thought-formatted instructional data from GSM8K and MATH, with initialization from DeepSeekMath-Instruct 7B. In spite of this narrow training focus, DeepSeekMath-RL 7B consistently yields better results than DeepSeekMath-Instruct 7B on all assessment criteria, underscoring the benefits imparted by reinforcement learning.

\section{Discussion}
This section details the main insights obtained from our pre-training and reinforcement learning investigations.

\subsection{Insights from Pre-Training}

We begin by summarizing our pre-training experiences. Unless mentioned otherwise, the experimental setup conforms to the standards specified in Section \ref{sec:quality-policy}. Specifically, references to the DeepSeekMath Corpus within this section pertain to the second data acquisition cycle, comprising 89B tokens.

\subsubsection{Impact of Code Training on Mathematical Reasoning}

A commonly held yet insufficiently validated belief posits that code training augments reasoning capabilities. Our study offers partial confirmation within the scope of mathematics: code training demonstrably enhances model proficiency in mathematical reasoning, regardless of tool integration.

To examine the influence of code-based training on mathematical reasoning, we designed experiments comparing two distinct setups: a two-stage training pipeline and a single-stage training approach.

\textbf{Two-Stage Training}

\begin{itemize}[topsep=0pt]
    \item \textbf{Code Training for 400B Tokens $\rightarrow$ Math Training for 150B Tokens}: DeepSeek-LLM 1.3B undergoes a first training phase on 400B code tokens, after which it is further trained with 150B tokens derived from math data.
    \item \textbf{General Training for 400B Tokens $\rightarrow$ Math Training for 150B Tokens}: For comparison, we include a baseline where the model is initially exposed to 400B tokens from a broad, general-purpose corpus curated by DeepSeek-AI rather than code-specific data. This setup seeks to clarify the distinct impact that code-centric pretraining imparts to mathematical reasoning, relative to generic text exposure.
\end{itemize}

\textbf{One-Stage Training}

\begin{itemize}[topsep=0pt]
    \item \textbf{Math Training for 150B Tokens}: The model is trained solely on 150B math tokens.
    \item \textbf{Training on a mixture of 400B Code Tokens and 150B Math Tokens}: A consecutive regimen—code training followed by math training—has been found to reduce coding ability. Here, we explore the effect of jointly presenting code and math tokens to the model during a single training period. This approach aims to determine whether concurrent exposure helps sustain both mathematical reasoning and coding capacities, while reducing catastrophic forgetting.
\end{itemize}

\paragraph{Results}
The quantitative outcomes from various training schemes are summarized in Table \ref{tab:code-math} and Table \ref{tab:code-math-general-eval}.

Across both staged and mixed training regimens, code-focused training enhances the model’s proficiency in program-assisted mathematical problem-solving. As indicated by the metrics in Table \ref{tab:code-math}, preliminary code training in the two-phase approach results in marked improvements on Python-based GSM8K and MATH benchmarks, with subsequent math-specific fine-tuning offering additional gains. Moreover, simultaneous exposure to code and math tokens during training (as in one-stage approaches) effectively counters catastrophic forgetting, which is a notable drawback of sequential training; this combined regimen also leads to enhanced coding (see Table \ref{tab:code-math-general-eval}) and programmatic math reasoning abilities (see Table \ref{tab:code-math}).

In addition, pretraining on code also supports non-tool-based mathematical reasoning. Within the two-stage framework, initial code token exposure yields noticeable advances even before math-specific instruction. It further amplifies the efficacy of math training in the subsequent phase, collectively achieving the strongest results. However, when code and math tokens are intermixed from the outset, the model's performance in pure mathematical reasoning without tools is diminished. A likely explanation is that the 1.3B parameter DeepSeek-LLM lacks adequate representational capacity to internalize both code and math domains in a simultaneous, blended training setting.

\subsubsection{ArXiv Papers Appear Unproductive for Advancing Mathematical Reasoning}

ArXiv papers have frequently been incorporated as part of the mathematical pre-training datasets \citep{minerva,gpt-f,llemma,mathpile}. Nevertheless, thorough investigations into their influence on mathematical reasoning capabilities are scarce. Interestingly, our experimental findings suggest that incorporating arXiv papers does not contribute positively to mathematical reasoning skills. To assess this, we utilized several model architectures, namely DeepSeek-LLM 1.3B and DeepSeek-Coder-Base-v1.5 7B \citep{deepseek-coder}, each trained on distinct arXiv datasets prepared using different preprocessing approaches:

\begin{itemize}[topsep=0pt]
    \item \textbf{MathPile} \citep{mathpile}: a dataset encompassing 8.9B tokens, assembled through rule-based cleaning and filtering, with scientific arXiv papers making up more than 85\% of the collection;
    \item \textbf{ArXiv-RedPajama} \citep{redpajama}: a set comprising all arXiv LaTeX source files, from which extraneous elements like preambles, comments, macros, and bibliographies have been removed, resulting in 28.0B tokens overall.
\end{itemize}

During experimentation, we separately pretrained DeepSeek-LLM 1.3B for 150B tokens and DeepSeek-Coder-Base-v1.5 7B for 40B tokens using each of these arXiv-focused corpora. Our findings consistently indicate that exposure solely to arXiv material fails to enhance, and sometimes even impairs, performance on a variety of mathematical evaluation benchmarks spanning multiple levels of difficulty. The assessed benchmarks encompass quantitative reasoning tasks (GSM8K and MATH, Table \ref{tab:arxiv-cot}), multiple-choice examinations (MMLU-STEM, Table \ref{tab:arxiv-cot}), and datasets involving formal mathematical reasoning (miniF2F, Table \ref{tab:arxiv_atp}).

Nevertheless, these observations come with caveats and must be interpreted cautiously. Several aspects remain unexamined, including:

\begin{itemize}[topsep=0pt]
    \item Potential advantages of arXiv-derived text for mathematical applications not represented in our analysis, such as the informalization of formal mathematical expressions and proofs;
    \item Interactions between arXiv data and additional training data of other genres;
    \item The emergence of any beneficial effects from arXiv content at scales exceeding those of our tested models.
\end{itemize}

Comprehensive investigation into these aspects remains a task for future research.

\subsection{Insights of Reinforcement Learning}
\subsubsection{Towards a Unified Framework} Here, we attempt to establish a generalized framework for analyzing training techniques, including SFT, RFT, DPO, PPO, and GRPO, and carry out experiments to study contributing elements within this unified view. Typically, for such training schemes, the gradient with respect to the parameters $\theta$ takes the following form:

\begin{equation}
    \nabla_{\theta}\mathcal{J}_{\textcolor{red}{\mathcal{A}}}(\theta) = \mathbb{E}[\underbrace{(q,o) \sim \textcolor{red}{\mathcal{D}}}_{Data \ Source}]\left( \frac{1}{|o|} \sum_{t=1}^{|o|}  \underbrace{GC_{{\mathcal{A}}}(q, o, t, \textcolor{red}{\pi_{{rf}}})}_{Gradient \ Coefficient}  \nabla_{\theta}\log \pi_{\theta}(o_t | q, o_{<t})\right).
\label{eq:objective}
\end{equation}

This framework involves three fundamental elements: 1) \textit{Data Source $\mathcal{D}$}, which specifies where training data is drawn from; 2) \textit{Reward Function $\pi_{{rf}}$}, responsible for providing feedback or reward signals during training; and 3) \textit{Algorithm $\mathcal{A}$}, which integrates both the data and reward signals, generating the gradient coefficient $GC$ that governs the degree of updating—either by penalizing or rewarding—each data instance. Within this unified framework, several core methods can be distinguished:

\begin{itemize}
    \item  \textbf{Supervised Fine-tuning (SFT)}: In SFT, a pretrained model is refined further by training it on SFT datasets curated by humans.
    \item \textbf{Rejection Sampling Fine-tuning (RFT)}: RFT extends SFT by continuing the fine-tuning process on model-generated responses that pass a filtering step based on the accuracy of answers to SFT dataset questions.
    \item \textbf{Direct Preference Optimization (DPO)}: DPO advances model performance by additional fine-tuning using augmented samples created from the SFT model, optimizing the model with a pairwise DPO loss function.
    \item \textbf{Online Rejection Sampling Fine-tuning (Online RFT)}: Unlike RFT, Online RFT initializes training from the SFT checkpoint but selects augmentation data in real time using the latest policy, refining the model with these contemporaneous samples.
    \item \textbf{PPO/GRPO}: PPO/GRPO employs the SFT-trained model as initialization and further improves it by reinforcing with outputs generated from the current version of the policy model.
\end{itemize}

The components underlying these approaches are outlined in Table \ref{tab:data_policy}. For a comprehensive explanation of the derivation steps, consult Appendix \ref{app:analysis-rl}.

\paragraph{Observation about Data Source} We classify data sources into two principal types: online sampling and offline sampling. Online sampling refers to collecting training examples through real-time interactions with the current training policy model, whereas offline sampling indicates the use of data obtained by sampling from the pretrained SFT model. The RFT and DPO frameworks utilize offline data, whereas Online RFT and GRPO are based on online data collection.

From Figure \ref{fig:rl-analysis}, it becomes apparent that Online RFT demonstrates marked superiority over RFT across both benchmarks. Initially, performance between Online RFT and RFT is quite similar; however, as training progresses, Online RFT establishes a clear lead, underscoring the efficacy of an online training regimen. This trend aligns with expectations—early in training, there is substantial similarity between the actor and the SFT model, with minimal difference in the samples collected. As training advances, the samples derived from the actor diverge more substantially, at which point the benefit of continual online data acquisition becomes more pronounced.

\paragraph{Observation about Gradient Coefficient} To update model parameters, the reward information is transformed into a gradient coefficient by the algorithm. Our analysis divides the reward mechanism into two forms: `Rule' and `Model'. The `Rule' method evaluates the appropriateness of an answer based on its factual accuracy, whereas the `Model' method employs a trained reward model to assign a score to each answer, with the reward model itself trained using data from the rule-based assessment. A central distinction between GRPO and Online RFT, as captured in Equations \ref{eq:GC-RFT} and \ref{eq:GC-GRPO}, is GRPO’s ability to adaptively modulate the gradient coefficient in response to the reward model’s output. This enables more granular encouragement or discouragement of answers according to reward magnitude. Conversely, Online RFT does not make such adjustments; incorrect answers receive no negative reinforcement, and all correct answers are equally reinforced regardless of their reward values.

As illustrated in Figure \ref{fig:rl-analysis}, GRPO demonstrates superior performance over online RFT, underscoring the effectiveness of modulating positive and negative gradient coefficients. Moreover, GRPO+PS outperforms GRPO+OS, which suggests that employing more granular, step-sensitive gradient coefficients yields additional advantages. We also examine the impact of iterative RL by running two consecutive iterations in our experimental setup. The results, visualized in Figure \ref{fig:iter-rl}, reveal that iterative RL delivers substantial performance gains, with the most pronounced improvements occurring after the initial iteration.

\subsubsection{Why RL Works?} In this work, we apply reinforcement learning to a subset of instruction tuning data, resulting in notable gains over models trained solely via instruction tuning. To elucidate the mechanism behind RL’s effectiveness, we assess both Pass@K and Maj@K metrics for Instruct and RL models on two benchmark datasets. According to Figure \ref{fig:maj-pass}, while RL leads to marked enhancements in Maj@K accuracy, Pass@K does not experience similar improvement. This pattern suggests that RL’s benefits primarily stem from making the output distribution more robust; specifically, \textbf{the improvement appears to derive from elevating the likelihood of correct responses among TopK predictions rather than fundamentally augmenting model competence.} Likewise, prior work such as \citep{wang2023making} identified a \textbf{misalignment problem} for reasoning tasks in SFT models, and demonstrated that preference alignment methodologies can be leveraged to bolster reasoning performance \citep{yuan2023rrhf,song2023preference,wang2023making}.

\subsubsection{How to Achieve More Effective RL?}
\label{sec:effec_rl}
Our experiments confirm that RL substantially benefits mathematical reasoning benchmarks. We also propose a general framework that unifies several representative training approaches, characterizing all as variants or simplifications of RL. This paradigm, encapsulated in Equation \ref{eq:objective}, centers on three essential elements: Data Source, Algorithm, and Reward Function. Below, we outline promising directions for each component.

\paragraph{Data Source} The data source constitutes the foundational input for any training paradigm. In RL settings, this specifically refers to unlabeled questions accompanied by responses sampled from the policy model. In this study, we restrict ourselves to instruction tuning questions and utilize a basic nucleus sampling strategy to generate responses. We hypothesize that this constraint may explain why our RL approach predominantly boosts Maj@K metrics. Moving forward, we aim to expand our RL framework to incorporate out-of-distribution prompts and integrate \textbf{advanced sampling (decoding) strategies}, such as those leveraging tree-search algorithms \citep{yao2023tree}. In addition, adopting \textbf{efficient inference techniques} \citep{xia-etal-2023-speculative,leviathan2023fast,kwon2023efficient,xia2024unlocking}—which are pivotal for enhancing exploration efficiency in policy models—represents another critical research avenue.

\paragraph{Algorithms} Algorithms leverage the available data and reward feedback to compute gradient coefficients, which in turn adjust the model parameters. As expressed in Equation \ref{eq:objective}, contemporary approaches largely \textbf{TRUST} the reward signal to guide the increase or decrease of a token’s conditional probability. Nevertheless, the assumption that the reward feedback is always accurate does not hold, particularly in highly complex domains. For instance, even rigorously labeled corpora such as PRM800K \citep{lightman2023let}—curated by experienced annotators—still exhibit roughly 20\% erroneous annotations\footnote{\url{https://github.com/openai/prm800k/issues/12\#issuecomment-1728491852}}. Consequently, it becomes necessary to investigate reinforcement learning methods capable of maintaining robustness to imperfect or noisy reward signals. We contend that alignment strategies characterized by a \textbf{WEAK-TO-STRONG} \citep{burns2023weak} framework will play a transformative role in the evolution of learning algorithms.

\paragraph{Reward Function} The reward function serves as the fundamental generator of training signals. In reinforcement learning, this is typically embodied by a neural reward model. We identify three crucial research avenues for reward models: 1) \textbf{Strategies for strengthening the reward model's generalization.} Robust generalization is vital for the reward model to address unseen tasks and novel decoding results; absent this, RL may only consolidate existing language model distributions rather than foster substantive improvements; 2) \textbf{Methods for quantifying and representing reward model uncertainty.} Such uncertainty could act as an essential mediator connecting suboptimal reward models to weak-to-strong learning frameworks; 3) \textbf{Approaches for rapid and effective construction of high-quality process reward models} that deliver detailed training feedback throughout reasoning chains \citep{lightman2023let,wang2023math}.

\section{Conclusion, Limitation, and Future Work} This work introduces DeepSeekMath, which establishes a new performance standard among open-source models on the MATH benchmark at the competition level, rivaling the results of proprietary counterparts. DeepSeekMath~builds upon the DeepSeek-Coder-v1.5 7B architecture and is continually trained for 500B tokens, incorporating 120B math-specific tokens harvested from Common Crawl as a substantial portion of its dataset. Our comprehensive ablation analysis demonstrates that web-based sources are rich in high-quality mathematical content, whereas data from arXiv falls short of anticipated utility. We propose Group Relative Policy Optimization (GRPO), an adaptation of Proximal Policy Optimization (PPO), which enhances mathematical reasoning abilities while reducing memory overhead. Empirical evaluations confirm that GRPO remains advantageous even for the high-scoring DeepSeekMath-Instruct 7B model. Additionally, we offer a cohesive framework for conceptualizing related approaches and identify promising research avenues for refining reinforcement learning methodologies.

While DeepSeekMath~excels in numerical reasoning tasks, its performance lags in domains such as geometry and theorem proving relative to leading proprietary systems. Notably, the model failed to solve questions involving triangles and ellipses in preliminary testing, which suggests that limitations in data curation during pre-training and fine-tuning may have introduced bias. Furthermore, the model’s relatively modest scale results in less robust few-shot learning compared to GPT-4, as DeepSeekMath~exhibits comparable effectiveness in zero-shot and few-shot scenarios, unlike the improvement seen in GPT-4 with few-shot examples. Our future work aims to further refine our data curation pipeline to assemble an even higher-quality pre-training corpus. We also plan to investigate directions outlined in Section \ref{sec:effec_rl} to develop more advanced reinforcement learning techniques for large language models.

\end{document}